\documentclass[10pt,a4paper]{article}
\usepackage[T1]{fontenc}
\usepackage[utf8]{inputenc}
\usepackage[english]{babel}
\usepackage{amsmath,amssymb,booktabs,geometry,hyperref}
\geometry{margin=1.5cm}
\hypersetup{colorlinks=true,urlcolor=blue,linkcolor=blue}

\title{Indexing Necessary Preconditions\\
for PDDL Successor Generation}
\author{Technical Report on \texttt{pddl\_fol}}
\date{September 24, 2026}

\begin{document}
\maketitle

\begin{abstract}
The \texttt{pddl\_fol} planner uses breadth-first search and evaluates every
grounded action in every state. We study an index that selects actions
using a necessary positive atom from the precondition, while leaving the
full first-order formula evaluation unchanged. On IPC 1998 Gripper 1,
median process time falls from 107.1 to 13.2 ms (8.14$\times$); on two
generated Gripper problems with 7 and 8 balls, the improvement reaches
12.13$\times$ and 13.80$\times$. Plans and explored-state counts match
the baseline. Results concern successor selection \emph{after} grounding;
they do not demonstrate a general improvement in FOL planning.
\end{abstract}

\section{Motivation and Related Work}
PDDL describes action schemas with variables and first-order formulas.
Many planners ground actions before search. Corr\^ea et al.~\cite{correa2020}
instead interpret lifted successor generation as query evaluation over a
state and study database optimizations. Lifted heuristics based on
delete relaxation and FF were proposed later~\cite{correa2021,correa2022}.
Fast Downward uses a representation and successor generators that differ
from this project~\cite{helmert2006}.

This study addresses a narrower bottleneck: the loop that scans every
already-grounded action at each expansion. It introduces no external FOL
solver, lifted search, or new heuristic. Treating applicability as a query
motivates the index, but the algorithm remains fully grounded.

\section{Index and Correctness}
Let $A$ be the sequence of grounded actions and $s$ the set of atoms
true in the state. For each action $a$, we seek a positive atom $r(a)$
that must be true whenever its precondition $\varphi_a$ is true. The
search traverses only atom and conjunction nodes; among candidates, it
chooses an atom of maximum arity. Subtrees containing disjunctions,
negations, implications, or quantifiers provide no constraints for this
index. An action with no necessary positive atom is placed in the
unfiltered set $U$.

Before search, we build $I_p=\{a\in A:r(a)=p\}$.
For each state, we evaluate only
\begin{equation}
 C(s)=U\cup\bigcup_{p\in s} I_p.
\end{equation}
We order $C(s)$ according to the original order in $A$, then evaluate
the full precondition for each candidate. If $a$ is applicable,
$\varphi_a(s)$ is true; by construction, $r(a)$ is true in $s$, or
$a\in U$. Therefore $a\in C(s)$: no applicable action is omitted.
Breadth-first search order and tie-breaking remain unchanged. The index
requires additional memory proportional to $|A|$; grounding cost is
unchanged.

\section{Experimental Protocol}
We compare initial commit \texttt{1487d5c} with the indexed branch,
both compiled using \texttt{cargo build --release} and Rust 1.92.0.
Machine: AMD Ryzen 5 7640U, Linux 7.2.5, 12 logical CPUs. Each measurement
uses \texttt{time.perf\_counter} around the full CLI process, including
startup, parsing, grounding, and search. Runs of the two versions are
interleaved; we report the median of 20 pairs for IPC instances and 3 pairs
for the two generated problems. Search limit: \texttt{--max-states 100000}.
CPU frequency was not controlled; short runtimes, especially Blocks, are
sensitive to noise.

The Gripper 1 (IPC 1998) and typed Blocks 4-0 (IPC 2000) domain and
problem files come from the public \texttt{pddl-instances} corpus,
revision \texttt{cf19edf7}~\cite{pddlinstances}. The repository files
\path{benchmarks/gripper-7.pddl} and \path{benchmarks/gripper-8.pddl}
extend the same domain to 7 and 8 balls; they are \emph{not} official IPC
instances. For every row, the baseline and indexed versions produce
identical CLI output, including the plan and number of explored states.

\section{Results}
\begin{table}[h]
\centering
\small
\begin{tabular}{lrrrrr}
\toprule
Problem & Baseline (s) & Index (s) & Speedup & States & Steps \\
\midrule
Blocks 4-0, IPC 2000 & 0.003592 & 0.002644 & 1.36 & 111 & 6 \\
Gripper 1, IPC 1998 & 0.107071 & 0.013155 & 8.14 & 253 & 11 \\
Gripper 7, generated & 5.344382 & 0.440760 & 12.13 & 4735 & 21 \\
Gripper 8, generated & 17.432608 & 1.262869 & 13.80 & 11773 & 23 \\
\bottomrule
\end{tabular}
\caption{Median end-to-end runtime on the same host. Speedup = Baseline/Index.}
\label{tab:ours}
\end{table}

The improvement grows on larger Gripper problems, where the old loop
evaluates many unnecessary action instances for every state. Blocks 4-0
is too small for firm conclusions from its 1.36$\times$ ratio. The
number of explored states remains unchanged: the index speeds up
successor generation; it does not reduce search.

\subsection{Comparison with Fast Downward}
We compiled Fast Downward 26.6+ from commit
\texttt{9b81c7e422fdf7be9f73b96e9a7a969c483dd5d4} in \texttt{release}
mode (GCC/G++ 16.2.1). We measure both \texttt{astar(blind())}, an
optimal reference without an informative heuristic, and the
\texttt{seq-sat-lama-2011} alias, a satisficing heuristic planner.
The \texttt{just install-fast-downward} command installs this revision
to \texttt{\~{}/.local/bin}; problems are identical to those in
Table~\ref{tab:ours}. FD times are medians of 20 runs for IPC cases and
5 for generated cases, on the same host; they include the Python driver,
translation, and search.

\begin{table}[h]
\centering
\small
\begin{tabular}{lrrr}
\toprule
Problem & Our index (s) & FD blind (s) & FD LAMA (s) \\
\midrule
Blocks 4-0, IPC 2000 & 0.002644 & 0.1324 & 0.1539 \\
Gripper 1, IPC 1998 & 0.013155 & 0.1267 & 0.1636 \\
Gripper 7, generated & 0.440760 & 0.1408 & 0.3196 \\
Gripper 8, generated & 1.262869 & 0.1402 & 0.6415 \\
\bottomrule
\end{tabular}
\caption{Median full-process runtime on the same host.
Fast Downward (FD) uses two different configurations.}
\label{tab:fd}
\end{table}

FD finds plans of the same length in every case; LAMA does not provide
a general optimality guarantee. On Gripper 8, FD blind is about 9 times
faster than our index in total runtime; LAMA is about 2 times faster.
Our executable has lower process latency on the two small IPC cases, but
startup and translation dominate FD: on Gripper 1, its blind search
takes only 0.000521 s, compared with 0.1267 s for the full process. This
comparison does not establish general superiority on these cases. With
blind, FD expands 11\,743 states on Gripper 8, compared with 11\,773 in
our search. These two FD configurations do not cover all its search
options.

\section{Limitations and Conclusion}
The evaluation covers two domains and a few problems on one machine.
The larger Gripper instances are generated; they do not replace a full IPC
suite or a timeout-based coverage comparison. The planner still generates
all parameter combinations before search, so the grounding problem
described by Corr\^ea et al.~\cite{correa2020} remains. Formulas without
necessary positive atoms use a full scan. The optimization offers a local,
measurable, and semantically conservative improvement; it does not
establish general competitiveness against mature planners.

\begin{thebibliography}{9}
\bibitem{correa2020}
A.~B. Corr\^ea, F. Pommerening, M. Helmert, and G. Franc\`es,
"Lifted Successor Generation Using Query Optimization Techniques,"
\emph{ICAPS}, vol. 30, pp. 80--89, 2020.
\url{https://doi.org/10.1609/icaps.v30i1.6648}.
\bibitem{correa2021}
A.~B. Corr\^ea, G. Franc\`es, F. Pommerening, and M. Helmert,
"Delete-Relaxation Heuristics for Lifted Classical Planning,"
\emph{ICAPS}, vol. 31, pp. 94--102, 2021.
\url{https://doi.org/10.1609/icaps.v31i1.15951}.
\bibitem{correa2022}
A.~B. Corr\^ea, F. Pommerening, M. Helmert, and G. Franc\`es,
"The FF Heuristic for Lifted Classical Planning,"
\emph{AAAI}, vol. 36, pp. 9716--9723, 2022.
\url{https://doi.org/10.1609/aaai.v36i9.21206}.
\bibitem{helmert2006}
M. Helmert, "The Fast Downward Planning System,"
\emph{Journal of Artificial Intelligence Research}, vol. 26,
pp. 191--246, 2006. \url{https://doi.org/10.1613/jair.1705}.
\bibitem{pddlinstances}
Potassco, \emph{PDDL Instances}, IPC 1998/2000,
\url{https://github.com/potassco/pddl-instances}.
\end{thebibliography}
\end{document}
